\cotefile{10}
\scanpage{123}
% archive sheet number: 122
% source lines 1--20
$\handwritten{\pi}_i^{\handwritten{2}}=\handwritten{\pi}_i,\handwritten{\pi}_i\handwritten{\pi}_j=0$ pour $i\handwritten{\ne}j$, $\handwritten{\sum_i\pi}_i=1$. La donnée d'une $\otimes$-structure sur $C$ impliquant les données a)b),et celles-ci compatibles avec le foncteur $C\handwritten{\to}D$ (pour $c=\mathrm{id}$),équivaut alors à une structure d'algèbre associative unitaire sur $A$, pompatible avec la structure de coalgèbre, et pour laquelle $A\handwritten{\to}k\handwritten{[}T,T^{-1}\handwritten{]}$ est un homomorphisme d'\unsure{algèbres|Àlgèbres}. Comment s'exprime le fait que \deleted{\unsure{$C\to D$ est}} sur $C$ il y a une structure de symmétrie \deleted{telle} respectée par le foncteur $C\handwritten{\to}D$ (et $c$ choisi plus haut) ? On peut d'abord exprimer la structure\inserted{\handwritten{$\otimes$}} précédente par une application diagonale de $U$ satisfaisant $\handwritten{\Delta}(\handwritten{\pi}_i)=\handwritten{\sum}_{j+k=i}\handwritten{\pi}_j\otimes\handwritten{\pi}_k$. \unsure{La|Ea} condition cherchée \deleted{revient à dire que si $M$ est un $U$-module de degré $i$, \unsure{N.e.|i.e.} tel que ${}_iM=M$, et $N$ de degré $j$ i.e. ${}_jN=N$, alors \unsure{làisomorphisme|laisomorphisme} $x\otimes y\quad(-1)^{ij}y\otimes x$ de $M\otimes_kN$ sur $N\otimes_kM$ est compatible avec les opérations de $U$} est celle-ci: si nous graduons $U$ par ``degré des opérateurs dans des gradués'' i.e. \deleted{$U^i\unsure{=\pi_j|=\sum_j}$} ces quotients $U_{\handwritten{\alpha}}$ par $U_{\handwritten{\alpha}}^i=\handwritten{\widehat{\sum}\pi}_{i+j}U_{\handwritten{\alpha}}\handwritten{\pi}_i$, alors l'application diagonale devient \uline{anti}commutative i.e. la \unsure{muttiplication|multiplication} de $A$ \deleted{\unsure{respecte}} est \uline{anticommutative}. (Il y aurait lieu de prouver que la graduation de $A$\inserted{i.e. de $U$} permet de récupérer le \deleted{\unsure{foncteur au|foncteur}} graduation sur le foncteur oubli,\margin{à vérifier, \handwritten{(\unsure{groupe} \unsure{de|des} classes \unsure{quelc.|quelq.} pour \unsure{graduation} ?)}} i.e. les $\handwritten{\pi}_i\handwritten{\in}U$ i.e. $A\handwritten{\to}k\handwritten{[}T,T^{-1}\handwritten{]}$, et que la graduation de $A$ \deleted{n'est} peut être choisie \unsure{arbitrairement}, avec la seule condition, de toutes façons nécessaire\deleted{s}, que $A$ soit anticommutative et l'homomorphisme diagonal \handwritten{\unsure{soit}} compatible aux graduations).

% source lines 21--27
\handwritten{9/} Quelles que soient les structures a)\deleted{\unsure{b)}}c) en jeu\inserted{(b) -l'unité - étant \handwritten{\unsure{fixée}}}, le groupe des $\otimes$-automorphismes du foncteur \deleted{fibre} ``oubli'' en est indépendant, c'est le groupe des $g\handwritten{\in}U$ tels que $\handwritten{\Delta}(g)=g\otimes g$ et d'augmentation 1\inserted{\handwritten{(i.e. les hom unitaires $\unsure{g|q}:A\to k$)}}\unsure{(}NB cette dernière condition exprime la condition b) de 3º). Remplaçant $k$ par un e algèbre quelconque, on trouve encore un \deleted{\unsure{foncte|foncteur}} schéma en monoïdes affine sur $k$, mais qui en général ne permettra pas de réconstitu\deleted{\unsure{t}}er la situation entièrement. \deleted{\unsure{(donner exemple}} Il faut analyser de plus près
